\section{引言}
\sourcepage{150}
到目前为止，我们主要研究“绝对”单纯同调群，尽管已经定义相对同调群并证明其拓扑不变性。现在将更详细地研究相对群。它们有许多用途：一方面，它们在拓扑的许多应用中自然出现；另一方面，要表述同调的基本性质，即 Eilenberg--Steenrod 公理，它们是不可缺少的。

相对同调群的一种用途，是提供关于绝对同调群的信息。$H_p(K,K_0)$ 与 $H_p(K),H_p(K_0)$ 之间存在联系。例如，若 $H_p(K,K_0)$ 和 $H_{p+1}(K,K_0)$ 都为零，则 $H_p(K)$ 与 $H_p(K_0)$ 同构。这并不显然，如何精确而一般地表述这种关系，是一个微妙的问题。

代数拓扑早期，这类定理的表述往往笨拙冗长，需要恰当的语言。Eilenberg 的一个出色代数思想使问题大为澄清：它实际上只是一种方便的新记号，称为群、环或其他代数结构的“正合列”。这一概念在代数与拓扑中的作用，怎么强调都不为过。含混的代数论证用正合列表述后往往变得十分透明；另一些连专业代数学家都觉得困难的论证，也变得足够直接，可以放心地留给读者。

若本章有关群的某些术语不熟悉，请参阅第 1.5 节。
